% ============================================================
% Twitch Sponsorship Optimization Agent -- Final Report
% MGMT 590-037 | AI-Enhanced Optimization
% Daniels School of Business, Purdue University
% ============================================================
\documentclass[11pt,letterpaper]{article}

% ── Packages ─────────────────────────────────────────────────────────────────
\usepackage[top=1in, bottom=1in, left=1in, right=1in]{geometry}
\usepackage[T1]{fontenc}
\usepackage[utf8]{inputenc}
\usepackage{mathptmx}
\usepackage{helvet}
\usepackage{booktabs}
\usepackage{array}
\usepackage{tabularx}
\usepackage{multirow}
\usepackage{xcolor}
\usepackage{amsmath}
\usepackage{fancyhdr}
\usepackage{titlesec}
\usepackage{enumitem}
\usepackage{parskip}
\usepackage{caption}
\usepackage{hyperref}
\usepackage{setspace}

% ── Colors ───────────────────────────────────────────────────────────────────
\definecolor{navy}{RGB}{0,32,91}
\definecolor{gold}{RGB}{207,163,71}
\definecolor{lightgray}{RGB}{245,245,245}
\definecolor{medgray}{RGB}{120,120,120}
\definecolor{darkgray}{RGB}{60,60,60}
\definecolor{flagred}{RGB}{192,0,0}

% ── Typography ────────────────────────────────────────────────────────────────
\setlength{\parskip}{4pt}
\setlength{\parindent}{0pt}

% ── Section formatting ────────────────────────────────────────────────────────
\titleformat{\section}
  {\large\bfseries\color{navy}}{\thesection}{0.75em}{}
  [\vspace{-6pt}\textcolor{gold}{\rule{\linewidth}{1.5pt}}\vspace{2pt}]

\titleformat{\subsection}
  {\normalsize\bfseries\color{navy}}{\thesubsection}{0.5em}{}

\titlespacing*{\section}{0pt}{10pt}{4pt}
\titlespacing*{\subsection}{0pt}{6pt}{2pt}

% ── Header and footer ────────────────────────────────────────────────────────
\pagestyle{fancy}
\fancyhf{}
\renewcommand{\headrulewidth}{0pt}
\fancyhead[L]{\small\textcolor{medgray}{Twitch Sponsorship Optimization Agent}}
\fancyhead[R]{\small\textcolor{medgray}{MGMT 590-037 \ $|$ \ Team 1}}
\fancyfoot[C]{\small\textcolor{medgray}{\thepage}}

% ── Table column types ───────────────────────────────────────────────────────
\newcolumntype{L}[1]{>{\raggedright\arraybackslash}p{#1}}
\newcolumntype{C}[1]{>{\centering\arraybackslash}p{#1}}
\newcolumntype{R}[1]{>{\raggedleft\arraybackslash}p{#1}}

% ── Caption style ────────────────────────────────────────────────────────────
\captionsetup{font=small, labelfont=bf, labelsep=period, skip=4pt}

% ── Hyperlink style ──────────────────────────────────────────────────────────
\hypersetup{colorlinks=true, linkcolor=navy, urlcolor=navy}

% ============================================================
\begin{document}

% ── Title block ──────────────────────────────────────────────────────────────
\thispagestyle{empty}

\begin{center}
  {\Large\bfseries\color{navy} Twitch Sponsorship Optimization Agent}\\[4pt]
  {\large\color{navy} Outcomes Report: Initial and Updated Problem Framing}\\[8pt]
  \textcolor{gold}{\rule{0.6\linewidth}{1.5pt}}\\[8pt]
  {\normalsize MGMT 590-037 \ $|$ \ AI-Enhanced Optimization \ $|$
   \ Daniels School of Business, Purdue University}\\[4pt]
  {\normalsize Team 1: Ethin DalPian \quad Nick Sopcic \quad Matt Foor \quad Shubh Vaishnav}\\[4pt]
  {\normalsize Summer 2026}
\end{center}

\vspace{6pt}
\textcolor{gold}{\rule{\linewidth}{1pt}}

% ── Executive Summary ────────────────────────────────────────────────────────
\section*{Executive Summary}
\vspace{-2pt}

We built an end-to-end AI agent that takes a sponsorship budget, fits audience
response curves to Twitch session data, and returns an optimal per-segment
allocation with a plain-language justification. Across both the initial framing
(V1) and the stakeholder-modified framing (V2), the agent consistently
outperforms intuitive heuristics.

\textbf{Key findings:}
\begin{itemize}[nosep, leftmargin=1.5em]
  \item At a \$20{,}000 budget the optimized agent produces \textbf{8.97
        incremental conversions}, a 2.3\% gain over historical allocation
        and a 630\% gain over a reach-chasing baseline, both statistically
        significant (bootstrap 95\% CI lower bound $>$ 0).
  \item The reach-chasing heuristic --- spend on the largest-follower
        segment first --- is the worst-performing strategy because the
        largest segment (Mega) has a poorly identified response curve and
        delivers near-zero measured lift.
  \item When per-segment conversion quality is incorporated (V2 stakeholder
        values: \$18--\$31 per conversion by tier), the revenue metric rises
        from \$23.53 to \textbf{\$249.99} at the same allocation, and the
        agent's revenue advantage over reach-chasing grows to \$222.94
        (median bootstrap uplift).
  \item A data-fusion estimate of per-segment values produced a near-flat
        vector (\$0.82--\$1.24) because the k-means clusters have similar
        session characteristics. The stakeholder-supplied tier values carry
        more information and are recommended for operational use.
\end{itemize}

\textbf{Primary recommendation:} deploy the agent with Model~B (stakeholder
per-segment values) and a budget of \$15{,}000--\$20{,}000, targeting
Segments~1 (Mid, \$31/conversion) and~3 (Micro, \$27/conversion). Expand
the creator pool in Segment~3 before increasing that segment's budget, as
its spend cap is already binding.

% ── Page 2 ───────────────────────────────────────────────────────────────────
\newpage

\section{Problem Framing (V1 --- Initial)}

\subsection{Business Problem}

A gaming studio wishes to allocate a fixed sponsorship budget across Twitch
creator segments to maximize incremental in-game conversions. Conversions are
game actions (purchases, sign-ups) that occur within three hours of a sponsored
session. The key challenge is that the relationship between sponsorship
spend and conversions is non-linear: audience response follows an S-shaped
Hill curve with increasing returns at low spend, an inflection point, and
diminishing returns at high spend. The resulting optimization is a non-convex
NLP, requiring a global search method rather than a standard gradient solver.

\subsection{Data and Segmentation}

The pipeline ingests 39{,}928 Twitch sessions across 2{,}000 creators and
12{,}000 player records. Creators are grouped into four segments using
k-means on log-transformed creator characteristics (followers, unique viewers,
session duration, sponsorship rate). The player dataset provides calibration-
period behavioral features used to estimate a per-conversion dollar value.

\begin{table}[h!]
\centering
\caption{Segment profile and fitted Hill curve parameters}
\small
\begin{tabular}{C{0.8cm} C{1.1cm} R{1.8cm} C{0.9cm} R{1.6cm} R{1.6cm} C{0.9cm} R{1.5cm}}
\toprule
\textbf{Seg.} & \textbf{Tier} & \textbf{Med.\ Followers} &
\textbf{$V$} & \textbf{$K$ (\$)} & \textbf{$n$} &
\multicolumn{2}{c}{\textbf{Spend Cap (\$)}} \\
\midrule
0 & Macro & 99{,}336 & 1.23 & 6.01 & 3.36 & \multicolumn{2}{c}{6{,}173} \\
1 & Mid   & 96{,}419 & 3.49 & 3.05 & 3.00 & \multicolumn{2}{c}{5{,}978} \\
2 & Mega  & 125{,}980 & 448.73 & 33{,}888 & 20.00$^\dagger$ & \multicolumn{2}{c}{14{,}612} \\
3 & Micro & 66{,}074 & 4.41 & 376.30 & 1.41 & \multicolumn{2}{c}{3{,}846} \\
\bottomrule
\end{tabular}\\[4pt]
{\small $^\dagger$\textcolor{flagred}{Segment~2 Hill fit hit the upper bound
($n = 20$). Parameters are poorly identified; results for this segment should be
treated as approximate.}}
\end{table}

\subsection{Optimization Formulation}

The agent solves a non-convex NLP:
\[
  \max_{\mathbf{x}} \sum_{s} g_s(x_s)
  \quad \text{s.t.} \quad \sum_s x_s \le B,\ 0 \le x_s \le u_s
\]
where $g_s(x_s) = V_s\, x_s^{n_s} / (K_s^{n_s} + x_s^{n_s})$ is the Hill
response curve for segment~$s$. With $n_s > 1$, each curve is S-shaped,
producing a combined objective surface with genuine local optima when budget
is split across segments in their convex regions. Basin-hopping with SLSQP
as the local minimizer is the primary solver; dual annealing serves as
validation. Both converged to identical objectives (gap $< 10^{-6}$), providing
strong evidence the global optimum was found.

The CVXPY library rejected the Hill objective as non-DCP, confirming
non-convexity. A single SLSQP run from a low-spend start was trapped at a
local boundary maximum (f\,=\,111.7), while basin-hopping escaped to the
global optimum (f\,=\,170.0), a 52\% improvement on the benchmark problem.

% ── Page 3 ───────────────────────────────────────────────────────────────────
\newpage

\section{V1 Results --- Initial Problem Framing}

\subsection{Prediction Layer}

The prediction pipeline (Component~2) estimates a per-session incremental lift
$\hat\tau$ using a T-learner: two separate XGBoost models, one trained on
sponsored sessions and one on unsponsored. The T-learner was selected over the
S-learner because it allows each treatment arm to capture arm-specific feature
interactions. Lift estimates from an IPW-reweighted S-learner confirmed
stability (pairwise correlation $\ge 0.83$ across all three estimators).

Hill saturation curves are then fitted to $(sp\_cost,\,\hat\tau)$ pairs per
segment via \texttt{scipy.optimize.curve\_fit}. A player value model
(two-stage XGBoost: classifier predicting spender status, regressor
predicting spend amount) yields a flat per-conversion value
$v = \$2.623$.

\subsection{Optimal Allocation}

\begin{table}[h!]
\centering
\caption{V1 optimal allocation at budget \$20{,}000 (flat $v = \$2.623$)}
\small
\begin{tabular}{C{0.8cm} C{1.1cm} R{1.7cm} R{2.0cm} R{2.0cm} C{1.4cm}}
\toprule
\textbf{Seg.} & \textbf{Tier} & \textbf{Spend (\$)} &
\textbf{Incr.\ Conversions} & \textbf{Incr.\ Revenue (\$)} &
\textbf{Funded} \\
\midrule
0 & Macro & 5{,}000.00 & 1.2295 &  3.23 & Yes \\
1 & Mid   & 5{,}000.00 & 3.4885 &  9.15 & Yes \\
2 & Mega  & 5{,}000.00 & 0.0000 &  0.00 & Yes$^\dagger$ \\
3 & Micro & 3{,}846.00 & 4.2517 & 11.15 & Yes \\
\midrule
\textbf{Total} & & \textbf{18{,}846.00} & \textbf{8.9697} & \textbf{23.53} & \\
\bottomrule
\end{tabular}\\[4pt]
{\small $^\dagger$\textcolor{flagred}{Segment~2 lift is effectively zero
due to the poorly identified Hill fit ($K = \$33{,}888$, cap $= \$14{,}612$).}}
\end{table}

Segment~3 hits its spend cap (\$3{,}846), the only binding constraint.
Segment~2 receives \$5{,}000 but delivers zero measurable lift, representing
an inefficient allocation driven by the unreliable Hill parameters.

\subsection{Baseline Comparison and Validation}

\begin{table}[h!]
\centering
\caption{Strategy comparison on held-out test creators (300 creators)}
\small
\begin{tabular}{L{3.4cm} R{2.0cm} R{2.4cm} R{2.3cm}}
\toprule
\textbf{Strategy} & \textbf{Spend (\$)} &
\textbf{Incr.\ Conversions} & \textbf{Incr.\ Revenue (\$)} \\
\midrule
Agent (Optimized) & 18{,}846.00 & 8.9697 & 23.53 \\
Historical        & 18{,}577.01 & 8.7711 & 23.01 \\
Reach-Chasing     & 20{,}000.00 & 1.2295 &  3.23 \\
\midrule
\multicolumn{2}{l}{Agent vs.\ Historical (\%)} & \textbf{+2.3\%} & \textbf{+2.3\%} \\
\multicolumn{2}{l}{Agent vs.\ Reach-Chasing (\%)} & \textbf{+629.9\%} & \textbf{+629.3\%} \\
\bottomrule
\end{tabular}
\end{table}

\begin{table}[h!]
\centering
\caption{Bootstrap uplift (1{,}000 resamples of held-out creators)}
\small
\begin{tabular}{L{3.0cm} R{2.2cm} R{2.2cm} R{2.2cm} C{1.2cm}}
\toprule
\textbf{Comparison} & \textbf{Median Uplift} & \textbf{95\% CI Lo} &
\textbf{95\% CI Hi} & \textbf{Win?} \\
\midrule
vs.\ Historical    & 0.1992 & 0.0959 & 0.3820 & \textcolor{navy}{\textbf{Yes}} \\
vs.\ Reach-Chasing & 7.7402 & 7.7402 & 7.7402 & \textcolor{navy}{\textbf{Yes}} \\
\bottomrule
\end{tabular}\\[4pt]
{\small Win declared when 95\% CI lower bound $> 0$.}
\end{table}

The agent wins against both baselines. The narrow confidence interval against
reach-chasing reflects deterministic segment follower rankings; regardless of
which creators appear in the bootstrap sample, the agent's avoidance of the
zero-lift Mega segment always produces a large advantage. The tighter margin
over the historical baseline is expected: the historical allocation already
approximates the optimum because it reflects practitioners' implicit knowledge
of which segments convert.

\textbf{Sensitivity highlights.} The objective is most sensitive to
Segment~3's saturation parameter~$V$ (max $|\Delta \text{obj}| = 0.425$
per $\pm 10\%$ perturbation). Segment~3's cap is binding; relaxing it by
50\% yields $+0.068$ additional conversions.

% ── Page 4 ───────────────────────────────────────────────────────────────────
\newpage

\section{Updated Problem Framing (V2 --- Stakeholder Modification)}

\subsection{Motivation}

The stakeholder identified that the flat value weight $v = \$2.623$ assigned
equal revenue importance to every conversion regardless of which creator
segment delivered it. Audience quality varies across tiers: niche, loyal
micro-creator communities tend to convert at higher lifetime value than
mega-creator audiences, who skew toward passive viewership. The modification
introduces a per-segment value vector $\{v_s\}$ and re-solves the allocation
under three value models.

\subsection{Phase 1 --- Stakeholder-Supplied Values}

Segments are ranked by median creator follower count and assigned to the
four stakeholder-defined tiers. The Mega tier --- despite the largest raw
audience --- receives the lowest per-conversion value (\$18), reflecting
lower audience spend propensity.

\begin{table}[h!]
\centering
\caption{Phase~1 segment-to-tier mapping and stakeholder values}
\small
\begin{tabular}{C{0.8cm} C{1.1cm} R{2.0cm} R{2.4cm} R{2.8cm}}
\toprule
\textbf{Seg.} & \textbf{Tier} & \textbf{Med.\ Followers} &
\textbf{Rank} & \textbf{$v_s$ -- Stakeholder (\$/conv.)} \\
\midrule
2 & Mega  & 125{,}980 & 1 (highest) & \$18.00 \\
0 & Macro &  99{,}336 & 2           & \$22.00 \\
1 & Mid   &  96{,}419 & 3           & \$31.00 \\
3 & Micro &  66{,}074 & 4 (lowest)  & \$27.00 \\
\bottomrule
\end{tabular}
\end{table}

\subsection{Phase 2 --- Data Fusion Estimate}

Because no shared key links the player and Twitch datasets, $v_s$ cannot be
estimated directly. We apply a bridging approach using \textbf{platform
engagement depth} as the shared construct: \texttt{pmins} (total gameplay
minutes) on the player side and \texttt{duration} (session length in minutes)
on the Twitch side. Both variables measure willingness to invest time in the
platform; neither is a follower-count or audience-size proxy.

The procedure ranks both variables on a $[0,1]$ percentile scale, fits a
binned mean $\hat{v}(\text{pmins\_percentile})$ on the player data, then
evaluates this function at the mean duration-percentile of each Twitch segment.

\begin{table}[h!]
\centering
\caption{Three-model value weight comparison}
\small
\begin{tabular}{L{2.8cm} R{2.0cm} R{2.0cm} R{2.0cm} R{2.0cm}}
\toprule
\textbf{Value Model} & \textbf{Seg~0 Macro} & \textbf{Seg~1 Mid} &
\textbf{Seg~2 Mega} & \textbf{Seg~3 Micro} \\
\midrule
(A) Flat $v$        & \$2.62 & \$2.62 & \$2.62 & \$2.62 \\
(B) Stakeholder $v_s$ & \$22.00 & \$31.00 & \$18.00 & \$27.00 \\
(C) Fusion $v_s$    & \$0.82 & \$1.24 & \$1.04 & \$0.95 \\
\bottomrule
\end{tabular}
\end{table}

\textbf{Fusion limitation.} The four k-means segments have similar mean session
durations (130--142 minutes), so the duration bridge produces a near-flat
fusion vector (range: \$0.43). A sensitivity check using a second bridge
(\texttt{chat\_rate\_pm} / viewer retention ratio) produced an equally flat
result. This confirms that the k-means segmentation captures engagement
quality dimensions that the available creator-level covariates do not
differentiate. Model~B (stakeholder values) carries materially more information.

\subsection{Three-Way Optimization Results}

The objective with per-segment values is $\max_{\mathbf{x}} \sum_s v_s\,g_s(x_s)$.
All three models were solved with basin-hopping ($n_\text{iter} = 200$) and
validated against dual annealing.

\begin{table}[h!]
\centering
\caption{Three-way comparison at budget \$20{,}000}
\small
\begin{tabular}{L{2.8cm} R{1.5cm} R{1.5cm} R{1.5cm} R{1.5cm} R{1.8cm} R{1.6cm}}
\toprule
\textbf{Model} & \textbf{Seg~0 (\$)} & \textbf{Seg~1 (\$)} &
\textbf{Seg~2 (\$)} & \textbf{Seg~3 (\$)} &
\textbf{Incr.\ Conv.} & \textbf{Revenue (\$)} \\
\midrule
(A) Flat     & 3{,}870 & 2{,}194 & 4{,}293 & 3{,}846 & 8.9697 & 23.53 \\
(B) Stakeholder & 3{,}870 & 2{,}194 & 4{,}293 & 3{,}846 & 8.9697 & 249.99 \\
(C) Fusion   & 3{,}870 & 2{,}194 & 4{,}293 & 3{,}846 & 8.9697 &   9.39 \\
\bottomrule
\end{tabular}
\end{table}

% ── Page 5 ───────────────────────────────────────────────────────────────────
\newpage

\section{V1 vs.\ V2: Analysis and Recommendations}

\subsection{V1 vs.\ V2 Comparison}

\begin{table}[h!]
\centering
\caption{Summary comparison: V1 initial framing vs.\ V2 updated framing}
\small
\begin{tabular}{L{3.2cm} L{5.4cm} L{5.4cm}}
\toprule
\textbf{Dimension} & \textbf{V1 --- Initial Framing} & \textbf{V2 --- Updated Framing} \\
\midrule
Value objective    & Flat scalar $v = \$2.623$ applied uniformly to all segments
                   & Per-segment vector $\{v_s\}$; three competing models (A/B/C) \\
\addlinespace
Optimization obj.  & $\max \sum_s g_s(x_s)$ (pure conversions)
                   & $\max \sum_s v_s\, g_s(x_s)$ (revenue-weighted) \\
\addlinespace
Allocation found   & Seg~0--2: \$5{,}000 each; Seg~3: \$3{,}846
                   & Seg~0: \$3{,}870; Seg~1: \$2{,}194; Seg~2: \$4{,}293; Seg~3: \$3{,}846 \\
\addlinespace
Revenue metric     & \$23.53 (flat $v$)
                   & \$249.99 (Model~B stakeholder $v_s$) \\
\addlinespace
Bootstrap uplift   & +0.199 conv.\ vs.\ Historical (WIN); +7.74 vs.\ Reach-chasing (WIN)
                   & +\$222.94 revenue vs.\ Reach-chasing for Model~B (WIN) \\
\addlinespace
Key limitation     & Hill fit for Seg~2 ($n = 20$) unreliable; Seg~2 assigned \$5{,}000 with zero lift
                   & Fusion estimate near-flat; stakeholder values preferred \\
\bottomrule
\end{tabular}
\end{table}

\subsection{Key Discussion Points}

\textbf{Q1: Are the flat-$v$ and pure-conversions allocations identical?}
Yes, mathematically and numerically. For any positive scalar $v$,
$\arg\max_{\mathbf{x}} v \sum_s g_s(x_s) = \arg\max_{\mathbf{x}} \sum_s g_s(x_s)$.
The two-norm difference between the V1 allocation (pure conversions) and
Model~A (flat-$v$ revenue) is \$0.00 to four decimal places. Model~A
therefore represents both objectives simultaneously.

\textbf{Q2: How do differentiated $v_s$ values move spend?}
At the tested budget (\$20{,}000), the physical allocation is unchanged across
all three models. This is a direct consequence of the near-saturated Hill
curves: with K values of \$3--\$376 for Segments~0, 1, and~3, all feasible
spend levels already operate in the diminishing-returns region, and the
optimizer cannot productively redirect capital because every segment's marginal
lift is similar. The revenue metric changes dramatically (from \$23.53 to
\$249.99) because the $v_s$ multipliers are 7--12$\times$ larger, not because
the allocation changed.

\textbf{Q3: Do Models~B and~C point in the same direction?}
Broadly yes: both assign the lowest weight to Segment~2 (Mega) and the
highest to Segment~1 (Mid). However, Model~C's range (\$0.43) is negligible
compared to Model~B's range (\$13), and both converge to identical allocations
in practice. Model~B is preferred because it encodes explicit audience-quality
knowledge derived from domain expertise. Model~C demonstrates the fusion
methodology but cannot recover meaningful differentiation from the available
creator-side covariates.

\subsection{Recommendations}

\begin{enumerate}[nosep, leftmargin=1.5em]

  \item \textbf{Deploy with Model~B (stakeholder values).} The revenue metric
  under stakeholder $v_s$ is the most informative objective. Use budget
  \$15{,}000--\$20{,}000; the agent's advantage over reach-chasing is robust
  across this range (bootstrap lower bound $>$ \$222).

  \item \textbf{Expand the Segment~3 creator pool.} Segment~3 (Micro,
  $v_s = \$27$, second-highest value) is the only binding cap constraint.
  Adding creators to this tier would allow the optimizer to allocate more
  capital to the most revenue-efficient audience without reducing spend
  elsewhere.

  \item \textbf{Re-fit the Segment~2 Hill curve with additional data.}
  The current fit is unreliable ($n = 20$, at the upper bound). The optimizer
  correctly assigns zero effective lift to Segment~2, but this means
  \$4{,}293--\$5{,}000 of each run's budget is either wasted or left unspent.
  Collecting more sponsored sessions with varied spend levels in Segment~2
  would allow a stable fit and potentially unlock material lift from the
  largest-audience tier.

  \item \textbf{Validate the lifetime-value assumption.} The flat value
  $v = \$2.623$ reflects first-transaction revenue only. Multiplying by a
  retention-adjusted lifetime-value factor (the agent stores
  $v_\text{adj} = \$1.03$ using observed churn rates) would reduce the
  absolute revenue estimates but not change the optimal allocation. Before
  budget decisions are made, the team should confirm which value definition
  the business uses for ROI calculations.

  \item \textbf{Accept the sponsorship-causality limitation.} Sponsorship
  is not randomized in this dataset. The T-learner and IPW robustness check
  control for observed confounders, but residual unmeasured confounding
  (e.g., seasonal creator popularity) means the lift estimates are
  observational, not causal. The agent's advantage over historical allocation
  is small (+2.3\%) precisely because the historical pattern already
  approximates a reasonable allocation. The large advantage over reach-chasing
  (\textgreater 600\%) reflects a structural insight---audience quality
  matters more than audience size---that is robust to confounding.

\end{enumerate}

\subsection{Conclusion}

The Twitch Sponsorship Optimization Agent demonstrates that a four-stage
pipeline (data, prediction, optimization, explanation) can convert raw session
logs into a defensible, revenue-aware sponsorship strategy. The core insight
running through both the V1 and V2 analyses is that the Mega-follower segment
is a trap: its large audience does not translate into measurable incremental
lift or high-quality conversions. Mid-tier and micro-niche creators deliver
more conversion value per dollar. Introducing per-segment value weights
(V2) sharpens this finding quantitatively---the revenue gap between the agent
and the reach-chasing heuristic grows from 630\% (conversions) to nearly
10$\times$ (revenue) under stakeholder values---and provides the business
with an actionable tier-differentiated budgeting framework.

\vspace{8pt}
\textcolor{gold}{\rule{\linewidth}{1pt}}
\vspace{4pt}

\noindent{\small \textbf{AI use disclosure:} This project utilized Claude
(Anthropic) for code generation, notebook scaffolding, and report drafting.
All model outputs, data interpretations, and recommendations were reviewed and
validated by the team.}

\end{document}
